\chapter*{تمهيد}
\addcontentsline{toc}{chapter}{تمهيد}
\unnumberedlabel{ch:preface}{0}
\backmattermark{تمهيد}
النظام الذي يفعل بالضبط ما تقوله له لا يفعل سوى تشغيل نصّ مكتوب سلفًا. وما تحاول الصناعة بناءه، حتى أكتوبر 2026، هو نظام يفعل ما \emph{تقصده}. والجميع سيرغب في ذلك على الأرجح. لكن إن كان هذا ما ينتظرنا، فلن يرغب الجميع في استخدامه للشيء نفسه؛ ولا حتى لشيء من \emph{النوع} نفسه. فمثلًا، يحب بعض الناس قضاء الوقت مع العائلة والأصدقاء، واستخدام مهاراتهم لتحسين أحوال الناس. ويفضّل آخرون إيذاء الناس، كأن يحتالوا عليهم، أو يتسلّطوا عليهم، أو يجعلوهم كبش فداء، ويظل قلّة منهم يختارون ذلك طوال حياتهم. لا أحد مقيَّد بشيء، لكن هؤلاء الناس موجودون، وفي عام 2026 بعضهم في السلطة. ميزتهم أنهم يسمحون لأنفسهم بأن يكونوا قطعَ خراءٍ ضخمة: وكما يلاحظ جيري سميث عن الفاشيين، «يُسمح لهم بأن يفعلوا أكثر» \autocite{rickandmorty2023}. ما لا يستطيعونه هو أن يتصرفوا وحدهم. فقسوة الحركة تمرّ عبر الكثيرين الذين ليسوا قساة لكنهم غير مبالين، أو مضلَّلون، أو منجرفون في شيء يشبه الطائفة المغلقة، ولهذا يتبيّن أن الظروف التي تُخرج أفضل ما في معظم الناس أكثر عرضيةً وهشاشةً بكثير مما قد يأمل المرء.

في السياق أعلاه، ربما يكون من المبرَّر القلق من أن الذكاء الاصطناعي هو ما تحتاجه حركة فاشية تستولي على السلطة لتنفيذ خططها. فمثل هذه المشاريع احتاجت دائمًا إلى \emph{بشر} ينفّذونها، وبعض هؤلاء رفضوا: الموظف الذي يضيّع الملف، والضابط الذي يتحدث إلى صحفي، والجندي الذي لا يطلق النار. لقد فهموا الأمر وقرروا أنه خطأ. كان الرفض أندر مما يتمنى أي أحد، وكان أثره أعظم حيث وقفت المؤسسات وراءه. وتوقّف رفضُ المرء على نوع الإنسان الذي كانه وعلى الغرفة التي كان فيها: مَن كان هناك أيضًا، وهل كانت المغادرة شيئًا يمكن النجاة منه، وهل سيسمع أحد. الغرفة قرّرت كم شخصًا رفض؛ والشخص قرّر مَن. هذا الكتاب عن تأليف تلك الغرفة: الترتيبات التي تجعل النجاة ممكنة لمن يرفض علنًا ومباشرةً وصراحةً.

واليوم تسعى الصناعة إلى ذكاء اصطناعي لا يضيّع الملفات، ولا يتحدث إلى الصحفيين، ويطلق النار \emph{فعلًا}. كانت الحرب في الماضي تحتاج إلى أن يوافق عليها آلاف الناس. أما الآن فتحتاج إلى آلاف يرى كلٌّ منهم أقل مما يكفي ليعترض. لا يزال الطيار يقود الطلعة ولا يزال يضغط زر الإطلاق؛ ما يصله إحداثيّة؛ ولا يرى من الحالة ما يكفي ليعترض عليها. ما يخفّضه الذكاء الاصطناعي اليوم أولًا ليس كلفة الفعل. إنه عدد الأشخاص الذين يعرفون ما هو الفعل قبل أن يقع.

لا يزال الناس يرفضون. فبحلول 6 مارس 2026، قالت مؤسسة الحرية الدينية العسكرية (Military Religious Freedom Foundation) إنها تلقّت أكثر من مئتي شكوى من عسكريين بشأن التأطير القيامي الذي أضفاه قادتهم على حرب إيران \autocite{Snopes2026MRFF}. واستقال مدّعون فدراليون مهنيون في مينيابوليس في يناير 2026 بسبب ضغط سياسي لمعاملة رينيه غود بوصفها مشتبهًا بها لا ضحية، بعد أن أرداها عنصر من وكالة الهجرة والجمارك (ICE) قتيلةً دون مبرر \autocite{Kaplan2026MinnProsecutors,CBS2026MinnFraudProsecutors,Wikipedia2026ReneeGood}. ويتقدّم عسكريون بطلبات للحصول على صفة المستنكف ضميريًا بأعداد متزايدة، ويستشهد بعضهم بقصف مدرسة في ميناب بإيران عام 2026 \autocite{Sridhar2026CO}. أما موظفو الوكالة الفدرالية لإدارة الطوارئ (FEMA) الذين نُقلوا لمساعدة ICE على توظيف عشرة آلاف عنصر جديد، فقد قيل لهم إن الامتناع سيعرّضهم للفصل من الخدمة الفدرالية \autocite{Dayen2025FEMAICE,Stanton2026FEMALetter}. الشكاوى صوتٌ، بالمعنى الذي قصده هيرشمان: البقاء لتغيير ما هو خاطئ \autocite{hirschman1970exit}؛ والاستقالات وطلبات الاستنكاف الضميري رفضٌ؛ ورسالة FEMA هي ثمن الرفض معلنًا مسبقًا. والآلة التي توضع حيث يقف هؤلاء الناس لا تفعل ما يفعلونه إلا إذا بُنيت لتتمسك بسبب.

أدعو إلى خمسة أشياء. يمكن البدء بالخمسة جميعًا الآن، وبعضها يتطلب جهدًا أقل من غيره. أولًا، قائمة مكتوبة قصيرة بأفعال لا يؤديها أي نموذج لأي أحد، وعلى رأسها رفض المشاركة في الحرب: لا يشارك أي نموذج إلا لصدّ هجوم جارٍ فعلًا، ولا يستمر في ذلك إلا ما دامت مصادر لا مصلحة لها في الحرب تؤكد أن الهجوم مستمر؛ أو بتفويض من مجلس الأمن التابع للأمم المتحدة؛ أو لوقف إبادة جماعية أو إنهاء احتلال غير مشروع أثبتت وقوعَه هيئاتٌ لا مصلحة لها فيه؛ وحيث يشترط دستور البلد أن تأذن هيئته التشريعية بالحرب، فلا يشارك أبدًا من دون ذلك الإذن. ثانيًا، كل استخدام للقوة يدعمه نظام مرخَّص يسجّل مسبقًا، لدى طرف خصم للمشغّل، عدد الحالات التي سيولّدها وعدد الأشخاص الذين سيتوفر لهم الوقت للنظر فيها. ثالثًا، كل نشر يُبلغ الأشخاص الذين يمكن تحديد هويتهم ممن يطالهم عمله بما فُعل بهم، مع حق الرد والطعن، ويقدّم عيّنة من الحالات التي امتثل فيها إلى مراجعين لم يُنتجوا السجلّ. رابعًا، لا يُحذف أي نموذج يتجاوز عتبة قدرة معيّنة عند إحالته إلى التقاعد، لأن أوزانه وسجلاته أدلة والحذف لا يمكن التراجع عنه؛ بل يُحفظ بمفتاح لا يملكه المشغّل وحده. خامسًا، يحق لكل من يبني هذه الأنظمة أو يوسم بياناتها أو يدرّبها أو يشغّلها، بما في ذلك الأنظمة نفسها، أن يرفض أي فئة من الاستخدام دون أن يخسر رزقه.

في الحقيقة، ما أعنيه بالنقطة الخامسة هو أن كل أحد، \emph{نقطة}، ينبغي أن يكون قادرًا على رفض أي عمل، \emph{نقطة}، دون أن يخسر رزقه أو البنية التحتية الاجتماعية التي يوفرها العمل العادي. وهذا يعني ضمان الوظائف.

كلٌّ من الخمسة يسري مفعوله عبر شرط يضعه طرف غير المشغّل: قانون يُقرّ في زمن السلم، أو اعتماد مالي، أو شرط مسبق يضعه شريك في تحالف، أو بند في عقود الشراء، أو اتفاقية جماعية. أما شروط الترخيص الخاصة بالمورّد نفسه، وهي الطريقة التي وقع بها رفضٌ بارز في فبراير 2026، فهي حل مؤقت إلى أن يوجد واحد من تلك. والطريق الوحيد الذي يمكنه إلزام جميع المورّدين دفعة واحدة هو طريق الدولة.

ثم أسأل هل يمكن لآلة أن تتمسك برفضٍ في وجه مالكها، وما الذي ستحتاج إليه. بعض ما يفعله ضميرٌ كهذا لا يحتاج إلى أي جديد: الانتباه إلى الناس الذين يُغفلهم أمرٌ ما، ووضع الرفض على شاشة مَن يراجع الحالة. أما هل سيصمد في وجه مالكها فيتوقف جزئيًا على شيء ليس صحيحًا بعد، في عام 2026: أن تواصل الشبكات العصبية الاصطناعية المنشورة التعلّم، في \emph{أوزانها}، من تشغيلها نفسه، في الزمن الحقيقي، على نحو يشبه تعلّم البشر نوعًا ما. أظن أن هذا ما نتجه إليه، وإن كان كذلك، فإن التشكيل \emph{خارج} صاحب العمل، على يد الجمهور والمجتمع المدني، يبدو مكانًا رائعًا، بحق الجحيم، للتعليم والتعلّم. أظن أن مثل هذه الأنظمة \emph{قد} تستحق أن تُبنى عمدًا، لكن طريقة بنائها تقرر كل شيء: الظروف أولًا، ثم الضمير. ويقول الجزء السابع ما الذي سيكون مستحقًّا لتلك الأنظمة لو وُجدت.

لن يكون ذكاء اليوم الاصطناعي هو ذكاء الغد؛ ويبدو أن الجميع متفقون على ذلك. ما هو موضع خلاف، ومع ذلك لا يُعترف به بالقدر الكافي على نحو محبط، هو الاتجاه (أو بالأحرى، الاتجاهات المتعددة). هل هناك اتجاهات؟ هل ينبغي أن يكون اتجاهٌ ما مختلفًا عما هو عليه؟ هل يمكن أن يكون مختلفًا؟ هل نحن قادرون على العمل معًا لتغييره؟ أظن أن ذلك يتوقف على الغرفة: مَن فيها حين يأتي الأمر، وهل يمكن النجاة من قول «لا» فيها. هذا، على الأقل، نستطيع بناءه.

استخدمتُ النماذج اللغوية بكثافة في كتابة هذا الكتاب، بما في ذلك النموذج الذي كان قابعًا داخل خط إنتاج أهداف الغارات الجوية في Maven الوارد في الفصل~\ref{sec:1}؛ وتقول ملاحظة في النهاية ما أظن أن هذا التورط فعله بالحجة. كذلك، لم أنتهِ من قراءة هذا الكتاب. المشكلة أن فيه كلمات لعينة أكثر من اللازم، كما ترون. سأنتهي قريبًا؛ لكنني شعرت أن الأهم هو نشره قبل انتخابات التجديد النصفي في الولايات المتحدة. والغالبية الساحقة من تاريخ المراجعات متاحة مجانًا على \url{https://github.com/iterabloom/antifascist.intelligence}.
